\documentclass[10pt,a4paper]{amsart}
\usepackage[margin=19mm]{geometry}
\usepackage{amsmath,amssymb,mathtools}
\usepackage[scr=rsfs,cal=euler]{mathalfa}
\usepackage{xfrac}
\usepackage[unicode,hidelinks,bookmarks=false]{hyperref}
\newcommand{\ie}{i.e.\ }
\newcommand{\cf}{cf.\ }
\newcommand{\resp}{respectively\ }
\newcommand{\Subsection}{\subsection}
\providecommand{\nobreakdash}{\nobreak\hbox{-}}
\setlength{\emergencystretch}{2em}
\multlinegap0pt
\usepackage{amsthm}
\theoremstyle{plain}
\newtheorem{theorem}{Theorem}
\newtheorem{proposition}{Proposition}
\theoremstyle{definition}
\newtheorem{definition}{Definition}
\theoremstyle{remark}
\newtheorem{remark}{Remark}
\newtheorem{notation}{Notation}
\DeclareMathOperator{\Pic}{Pic}
\DeclareMathOperator{\rk}{rk}
\DeclareMathOperator{\irreg}{irreg}
\DeclareMathOperator{\reg}{reg}
\newcommand{\ce}{\mathrel{\mathop:}=}
\title[Irregular bundles with prescribed weights on Hopf surfaces]{Irregular bundles on Hopf surfaces: the existence of irregular bundles with prescribed weights (Theorem 4)}
\author{Edoardo Ballico, Elizabeth Gasparim}
\date{Source: arXiv:2508.21814v1 (CC BY 4.0)}
\begin{document}
\begin{abstract}
We discuss the hypersurfaces of the moduli spaces of rank $2$ vector bundles 
on a classical Hopf surface formed by irregular bundles.


\end{abstract}
\maketitle

\noindent\emph{Source and licence.} Irregular bundles on Hopf surfaces. Version arXiv:2508.21814v1 (CC BY 4.0), distributed under the Creative Commons Attribution 4.0 International licence, http://creativecommons.org/licenses/by/4.0/, as recorded in the arXiv licence metadata for this exact version and shown on the abstract page https://arxiv.org/abs/2508.21814v1. Full rights notice: licenses/arxiv-2508.21814-CC-BY-4.0.txt.\par\smallskip

\noindent\textit{Editorial scope.} Counted unit: the theorem of the article printed as Theorem 4 in the arXiv version and announced as item (4) in its introduction (source0.tex lines 652--659, source label Wrr7), delivered together with the article's complete original proof of it (source0.tex lines 661--716, one \texttt{proof} environment of 56 lines that proves both part (a) and part (b) with no gap and no deferral). No mathematics is written, repaired or re-derived: statement and proof are the article's own text line for line, in the article's own notation ($n$, $\mathcal M_n(0)$, $\mathcal M_n^{\irreg}(0)$, $|\mathcal O_{\mathbb P^1\times \mathbb P^1}(n,1)|$, $|_{sm}$, $G$, $D_i$, $w_i(E)$, $w(E)$, $\Gamma$, $\Delta$, $\mathcal M_n(0)$). Required context retained uncounted: the article's abstract (lines 146--147); its moduli space $\mathcal M_n$ with the spectral involution, the double cover (line 313--342, including the definition of $\mathcal M_n(0)$ as the bundles without jumps); its definition of $\mathcal M_n^{\reg}(0)$ (lines 344--347) and of $\mathcal M_n^{\irreg}(0)$ together with the three identities identifying those sets with the inverse images under the graph map (lines 357--362); its remark identifying $G(\mathcal M_n(0))$ with the Zariski open set of smooth graphs (lines 350--352); the opening of its section on irregular profiles, with the four thetas, the four points $a_1,\dots,a_4$, the divisors $D_i=\pi_2^{-1}(a_i)$ and the set $|\mathcal O_{\mathbb P^1\times \mathbb P^1}(n,1)|_{sm}$ of smooth elements (lines 461--469); the four zero-dimensional schemes $Z_U(i)=U\cap D_i$ (lines 471--473) and the remark that a bundle is irregular exactly when one of them fails to be reduced (lines 475--477); the notation fixing the tangency sets $H_i$ and their inverse images $\mathcal H_i$ (lines 479--482); the article's definition of the $i$-th profile, the $i$-th weight, the length and the reduced profile (lines 497--511) together with the sentence relating the weight to the reduced profile (lines 513--515); the notation fixing the multiplicities and the weights $w_i(E)$ and $w(E)$ (lines 517--523); the statement of the article's theorem printed as Theorem 3 (lines 568--578), which the counted proof invokes for the case $n=2$ and for the bound $w(E)\le 2n-2$; the statement of the article's proposition printed as Proposition 1 (lines 639--641), which the counted proof invokes by name; and the article's section heading ``Ramifications and weights'' (line 635), under which the counted theorem is printed. Every definition, numbered display, label and structural statement of those lines is retained unchanged. Omitted from this package and not redistributed: the article's front matter (its titlepage block, addresses and e-mail addresses, lines 126--145), its keywords line and table of contents, its first three sections with the whole of its introduction and of its sections on bundles without jumps and on regular versus irregular bundles (lines 148--312 and 363--460, which contain the article's other announced theorems, its spectral-curve section and its regular-versus-irregular section), the parts of its fourth and fifth sections that lie between the retained excerpts (lines 148--312, 343, 348--349, 353--356, 363--460, 470, 474, 478, 483--496, 512, 516, 524--567, 579--634, 636--638, 642--651, 660), the whole of its sixth and seventh sections (lines 717--877), its acknowledgements and its bibliography (lines 878--982), and its closing \texttt{\textbackslash end\{document\}}. Nothing inside a retained excerpt is omitted or altered; the unit's own text outside the excerpts is the wrapper matter described below. Citations: the retained excerpts carry seven citation commands, of which six are active and cite exactly three keys of the article's own bibliography, \texttt{BH}, \texttt{BMo} and \texttt{h} (the seventh sits inside the article's own commented-out line about the spectral involution and stays a comment in the delivery). Those three entries are transcribed verbatim from the article's own in-source \texttt{thebibliography} (source0.tex lines 879--982, 27 entries) in the article's own order and with the article's own printed labels, so the delivered citation marks read as the article's do (\texttt{[BH]}, \texttt{[BMo]} and \texttt{[H]}), and the entry for \texttt{h} is delivered with the article's own printed misspelling of the author's name. No other entry of the article's bibliography is transcribed or redistributed. Reference adaptation: none was needed and none was made. Every reference inside the delivered unit resolves inside it: the label \texttt{Wrr6} of the retained Proposition 1, which the counted proof cites by name; the label \texttt{Wrr5} of the retained Theorem 3, which the counted proof cites by name; the label \texttt{dc} of the retained double-cover display, which the retained sentence about the four points $a_1,\dots,a_4$ cites; and the label \texttt{eqb2} of the numbered display inside the counted proof, which the proof cites twice later on. The labels \texttt{defH}, \texttt{weights} and \texttt{eqb1} are retained as the article's own anchors and are not cited anywhere in the unit, which produces no unresolved reference. No unresolved reference or citation is produced. Printed slips kept literal: no printed word, symbol, hypothesis or number of the counted statement or of its proof was altered. In particular the counted proof calls the retained Proposition 1 a lemma (``By Lemma \texttt{\textbackslash ref\{Wrr6\}}''), exactly as the article prints it; the proof's phrase ``Since the ramification points of $X$ have different,'' is delivered with the article's own missing complement; the article's own spelling \texttt{Harshorne} in the entry for \texttt{h} is kept; the article's own uses of $\ce$ for the definition sign and of the vertical bar in restrictions are kept; and the article's own mislabelled cross-reference ``Lemma \texttt{\textbackslash ref\{Wrr6\}}'' prints the unit's own number for the retained Proposition 1. Layout and class: the article's own class is \texttt{amsart} with the options \texttt{a4paper}, \texttt{reqno}, \texttt{oneside} and \texttt{11pt}, and its own preamble loads, among others, \texttt{xcolor}, \texttt{tikz}, \texttt{tikz-cd}, \texttt{lmodern}, \texttt{braket}, \texttt{microtype}, \texttt{hyperref} and \texttt{cleveref}; the wrapper is the lane's pinned \texttt{amsart} editorial wrapper at 10pt on A4, and no class or style file travels with this package because the article uses none beyond \texttt{amsart} itself. The wrapper adds the article's own macro definitions that the retained text needs (the operators $\Pic$, $\rk$ and $\irreg$ and the definition sign $\ce$), declares the theorem environments the retained statements need, and reproduces the article's \texttt{amsthm} environment set. The article's custom theorem style (a small-caps head with no body font change) is not reproduced, so the retained environments print in the class's standard styles, and the delivered numbering is the unit's own: the retained Theorem 3 prints as Theorem 1 of the unit, the counted Theorem 4 as Theorem 2, and the retained Proposition 1 as Proposition 1. The article's font setup (\texttt{lmodern}) coincides with the wrapper's Latin Modern text and math fonts, and its 11pt measure differs from the wrapper's 10pt; that is a page-appearance difference of the wrapper only, and no formula, symbol, hypothesis or argument changes. Assets: the article has no figure environment, no graphics-inclusion command, no PDF-page inclusion, no table and no TikZ picture anywhere in its source (its statements are purely verbal and symbolic); no retained span contains a graphics-inclusion command, an \texttt{input} directive or an \texttt{include} directive; the package declares no asset dependency and no image file is copied into or redistributed with it. All mathematics of the unit is native text with no image dependency. Licence: version arXiv:2508.21814v1 of this article is distributed under the Creative Commons Attribution 4.0 International licence, http://creativecommons.org/licenses/by/4.0/, as recorded in the arXiv licence metadata for this exact version and shown on the abstract page https://arxiv.org/abs/2508.21814v1. A full rights notice is delivered at licenses/arxiv-2508.21814-CC-BY-4.0.txt. Characters: the retained excerpts contain exactly one character outside ASCII, the accented letter of the word Poincar\'e as the article's own source writes it in its construction of the graph map; the delivered engine is XeLaTeX, whose default Latin Modern text font carries that character, and the compile receipt reports no missing character. The article's bibliography escapes its other accents with ASCII control sequences, and those transcriptions are delivered as the article writes them. No glyph is substituted and no word is rewritten.
Let $\mathcal M_n$ denote the moduli space
of all rank $2$ stable vector bundles on $X$ with trivial determinant and $c_2(E)=n$, i.e.
$$\mathcal M_n=\{E\,  \text{stable  bundle on} \, X,   \rk(E)=2, \det(E)\cong \mathcal O_X, c_2(E)=n\}.$$




The variety $\mathcal M_n$ is  smooth non-empty    of dimension 
$4n$ \cite[Prop.\thinspace 3.4.4]{BH}, see also 
\cite[Prop.\thinspace 4.2]{BMo} for the case of nontrivial determinant.
%A spectral involution on  is defined  defined 
%even if $\delta \ne \mathcal O_X$ \cite[p.\thinspace 3]{BMo}.
Since here we consider only the case of trivial determinant,  the 
{\bf spectral involution} on $T^*$ may be defined 
as $i(\lambda) = -\lambda$, with corresponding {\bf double cover} 
\begin{equation}\label{dc}\mathbb P^1 \times T^* \to \mathbb P^1 \times \mathbb P^1\end{equation} 
$$(b, \lambda) \stackrel{\sigma}{\mapsto} (b, \{\lambda, i(\lambda)\}).$$


Given a rank 2 bundle $E$ on $X$, let $V$ be the universal Poincaré line bundle on
$X \times \mathbb C^*$ and consider the first derived image $R^1\pi'_*(E \otimes V)$ where
$\pi'$  the projection $\pi'= \pi\times id \colon X\times \mathbb C^* \to \mathbb P^1 \times \mathbb C^*$. 
The sheaf $R^1\pi'_*(E \otimes V)$ is supported on a divisor on $ \mathbb P^1\times \mathbb C^*$, 
which descends to a divisor 
$$D \subset \mathbb P^1\times \mathbb P^1 = \pi(X) \times \Pic^0(T)/\pm1.$$ When $c_2(E)=n$, the divisor
 $G(E) \ce D $ belongs to the linear system $|\mathcal O(n,1)|$ over $\mathbb P^1\times \mathbb P^1$
 and is called the {\bf graph} of $E$, see\cite[Prop.\thinspace 3.2.3]{BH}. \\

Let  $\mathcal M_n(0)$ denote the set of all $E\in  \mathcal M_n$ with {\bf no jumps}, i.e.
$$\mathcal M_n(0) \ce  \{E\in  \mathcal M_n,  \forall T_x \, E\vert_{T_x}\, \text{ is semistable}\}.$$

 Let  $\mathcal M_n^{\reg}(0)$ denote the subset of all {\bf regular bundles} in $\mathcal M_n(0)$, that is, 
 those whose restrictions to fibers are semistable but not sums of  half-periods, i.e.
 $$\mathcal M_n^{\reg}(0)= \{E\in  \mathcal M_n: \forall   T_x, 
  E\vert_{T_x}   \not\cong  \, L_0\oplus i^\ast(L_0) \, \text{with}\, L_0^{\otimes 2} \cong \mathcal O_{T_x} \}.$$

\begin{remark}The image $G(\mathcal M_n(0))$ of $\mathcal M_n(0)$ is the Zariski open set of all graphs $D\in |\mathcal O_{\mathbb P^1\times \mathbb P^1}(n,1)|$ which are smooth (or, equivalently in this case, irreducible), i.e. the set of all $D\in |\mathcal O_{\mathbb P^1\times \mathbb P^1}(n,1)|$ 
giving a degree $n$ morphism $\mathbb P^1\to \mathbb P^1$.
\end{remark}

Assume $n>1$ and set $$\mathcal M_n^{\irreg}(0)\ce \mathcal M_n(0)\setminus \mathcal M^{\reg}_{\delta,n}(0).$$ 
Since the notions of irregularity and of having no jumps are completely 
 described in terms of graphs, we have the equalities
$$ \mathcal M_n(0) =G^{-1}(G(\mathcal M_n(0))), \quad \quad 
 \mathcal M_n^{\reg}(0) =G^{-1}(G(\mathcal M_n^{\reg}(0))),  $$and
$$ \mathcal M_n^{\irreg}(0) =G^{-1}(G(\mathcal M_n^{\irreg}(0))).$$

\section{Irregular profiles}
We fix $\delta\in \mathrm{Pic}(X) =\mathrm{Pic}^0(X)\cong \mathbb C^\ast$. Let $I\subset T^\ast$ be the set of all thetas with respect to the involution $i_\delta$.
We we fix an order for the $4$ elements of $I$,  we write them as $L_1,L_2,L_3,L_4$, and write $a_1, a_2, a_3, a_4$ 
for the $4$ elements of $\mathbb P^1$ associated to them, see \eqref{dc}.
Let $\pi_2\colon \mathbb P^1\times \mathbb P^1\to \mathbb P^1$ denote the projection onto the second factor. 
Set 
$$D_i\ce \pi_2^{-1}(a_i).$$ We get $4$ elements of $|\mathcal O_{\mathbb P^1}(1,0)|$. 

Let  $|\mathcal O_{\mathbb P^1\times \mathbb P^1}(n,1)|_{sm}$ denote the set of smooth elements in $|\mathcal O_{\mathbb P^1\times \mathbb P^1}(n,1)|$.

 Note that $E\in \mathcal M_n(0)$, i.e. $E$ has no jumps, if and only if $G(E)\in |\mathcal O_{\mathbb P^1\times \mathbb P^1}(n,1)|_{sm}$. For each $U\in |\mathcal O_{\mathbb P^1\times \mathbb P^1}(n,1)|_{sm}$ we get $4$ degree $n$ zero-dimensional schemes
 $$Z_U(i)\ce U\cap D_i, \quad  i=1,2,3,4.$$
  For any $E\in \mathcal M_n(0)$, taking $G(E)$ in place of $U$,  we get $4$ zero-dimensional schemes $Z_{G(E)}(i)$, $i=1,2,3,4$.

  \begin{remark}
Note that $E$ is irregular if and only if at least one of the 4 schemes $Z_{G(E)}(i)$, $i=1,2,3,4$, is not reduced, i.e., it is not formed by $n$ distinct points.
\end{remark}

\begin{notation}\label{defH} For each $i=1,2,3,4$, denote by  $H_i$  the set of all divisors 
$D\in |\mathcal O_{\mathbb P^1\times \mathbb P^1}(n,1)|_{sm}$ such that $D$ is tangent to $D_i$ 
(the point of tangency with $D_i$ is not fixed), and let $\mathcal H_i\ce G^{-1}(H_i)$.
\end{notation}

\begin{definition}\label{weights}
Let $m_1\ge \cdots \ge m_x>0$ be the degrees  of the connected components of $D\cap D_i$
(in decreasing order).
\begin{itemize}
\item  We say that $(m_1,\dots ,m_x)$ is the {\bf $i$-th    profile} of $D$ and of all bundles $E$ with $G(E)=D$. 
 
\item  We say that the integer $\sum (m_i-1) =m_1+\cdots +m_x-x$ is the {\bf $i$-th weight} of $D$ and of all bundles $E$ with $G(E) =D$. 
 
\item  The {\bf length} $\ell$ of the profile $D\cap D_i$ and of all  bundles $E$ with $G(E)=D$
is the maximal integer $y$ such that $m_y\ge 2$, with the convention $\ell =0$ if $m_1=1$, i.e. if $D$ is transversal to $D_i$. 

\item The {\bf reduced $i$-th profile} of $D$
is the set of integers $m_1,\dots ,m_\ell$.
\end{itemize}
\end{definition}

 If $\ell>0$, these numbers form a non-decreasing sequence of integers $\ge 2$ with $\sum _i m_i\le n$. 
 Note that the $i$-th weight is uniquely determined by the reduced profile of $m_1,\dots ,m_s$: it is $0$ if $\ell =0$, while it is $-\ell+m_1+\cdots +m_\ell$ if $\ell>0$.
For instance, for $n=2$ the possible profiles are $2$ and $1,1$; while for $n=3$ the profiles are $3$, $2,1$ and $1,1,1$.\\

\begin{notation}For any such $E$ with no jump we denote by 
$$m_1,\dots ,m_{s(i)} \quad  \text{or} \quad  m_1(E),\dots ,m_{s(i)}(E)$$
 the multiplicities of the connected components of $Z_{G(E)}(i)$ in non-decreasing order. 
By definition, the $i$th-weight $w_i(E)$ of $E\in \mathcal M_n(0)$ is the integer 
 $$w_i(E)\ce \sum _{j=1}^{s(i)} (m_j-1),$$
 and the weight $w(E)$ of $E\in \mathcal M_n(0)$ is the integer $$w(E)\ce w_1(E)+w_2(E)+w_3(E)+w_4(E).$$
\end{notation}

\begin{theorem}\label{Wrr5}
Assume $n>1$. For all $E\in \mathcal M_n^{\irreg}(0)$  we have:
\begin{enumerate}
\item $w(E)\le 2n-2$.
\item $E$ has $(n, 0,\dots, 0)$ has its $i$-th profile if and only if $w_i(E)\le n-1$ for all $i$ and $w_i(E)=n-1$ .
\item Fix $i\in \{1,2,3,4\}$, a positive integer $\ell$ and a restricted profile $(m_1,\dots ,m_\ell)$. 
There exists a non-empty locally closed analytic subset $\Gamma \in \mathcal M_n^{\irreg}(0)$
such that  all  $G(E)\in |\mathcal O_{\mathbb P^1\times \mathbb P^1}(n,1)|$ have   $(m_1,\dots ,m_\ell)$ as $i$-th profile
and this set of graphs $G(E)$  has codimension at most $-\ell +m_1+\cdots +m_\ell$ in $|\mathcal O_{\mathbb P^1\times \mathbb P^1}(n,1)|$.
\end{enumerate}
\end{theorem}

\section{Ramifications and   weights}

\begin{proposition}\label{Wrr6}
Fix an integer $n\ge 2$. Take a general $D\in |\mathcal O_{\mathbb P^1\times \mathbb P^1}(n,1)|$. Then $\pi_{2|D}: D\to \mathbb P^1$ has $2n-2$ distinct ramification points and their images are $2n-2$ distinct points of $\mathbb P^1$.
\end{proposition}

\begin{theorem}\label{Wrr7}
Fix an ordering $i_1,i_2,i_3,i_4$ of the set $\{1,2,3,4\}$.
\begin{itemize}
\item[(a)] For all $n\ge 2$ there exists  a codimension $2$ locally closed analytic subset $\Gamma \subset \mathcal M_n(0)$ such that $w_{i_1}(E)=w_{i_2}(E)=1$ and $w_{i_3}(E)=w_{i_4}(E)=0$ for all $E\in \Gamma$.

\item[(b)] For all $n\ge 4$ there exists  a codimension $3$ locally closed analytic subset $\Delta \subset \mathcal M_n(0)$ such that $w_{i_1}(E)=w_{i_2}(E)=w_{i_3}(E)=1$ and $w_{i_4}(E)=0$ for all $E\in \Delta$.
\end{itemize}
\end{theorem}

\begin{proof}
To simplify the notation we use $j$ instead of $i_j$. The other permutations only need more double or triple indices.

We always work in $|\mathcal O_{\mathbb P^1\times \mathbb P^1}(n,1)|_{sm}$. The surjectivity of the graph map $G$ \cite[Th. 5.2.2]{BH} gives that it sufficient to find
$\Gamma_1\subset |\mathcal O_{\mathbb P^1\times \mathbb P^1}(n,1)|$, $n\ge 2$, (resp. $\Delta_1\subset |\mathcal O_{\mathbb P^1\times \mathbb P^1}(n,1)|$, $n\ge 4$) such that
$\dim \Gamma _1 =\dim |\mathcal O_{\mathbb P^1\times \mathbb P^1}(n,1)| -2$ (resp. $\dim \Delta _1 =\dim |\mathcal O_{\mathbb P^1\times \mathbb P^1}(n,1)| -3$) such that all $C\in \Gamma_1$ (resp. $C\in \Delta_1$) have this
 property.
By Lemma \ref{Wrr6} there is $A\in |\mathcal O_{\mathbb P^1\times \mathbb P^1}(n,1)|_{sm}$ such that $\pi_{2|A}$ has $2n-2$ ramification points with $2n-2$ different images $a_1,\dots ,a_{2n-2}\in \mathbb P^1$. The group $\mathrm{Aut}(\mathbb P^1)$ acts uniquely 3-transitively on $\mathbb P^1$, i.e. for any $2$ triples $(b_1,b_2,b_3)$, $(c_1,c_2,c_3)$ of distinct points of $\mathbb P^1$ there is a unique $\gamma\in \mathrm{Aut}(\mathbb P^1)$ such that $\gamma(b_j)=c_j$ for all $j$.\\

 (a) In this step we prove part (a).  We saw that
there is $\alpha_2 \in \mathrm{Aut}(\mathbb P^1)$ such that $\alpha_2(a_1) =\pi_2(D_1)$ and $\alpha_2(a_2)= \pi_2(D_2)$. Let $\alpha$ denote the automorphism of $\mathbb P^1\times \mathbb P^1$ which acts as the identity 
on the first factor and as $\alpha_2$ on the second factor. For all $E\in \mathcal M_n$ such that $G(E)=\alpha(A)$
we have $w_1(E)=w_2(E)=1$, because all ramification points of $D:= \alpha(A)$ are ordinary, 
exactly one of  them is contained in $D_1$ and exactly one of them is contained in $D_2$. 
If $n=2$ we also have $w_3(E)=w_4(E)=0$, because $w(E)\le 2$ for all $E\in \mathcal M_{2}(0)$ 
(Theorem \ref{Wrr5}). However, for $n>2$ to get $w_3(E)=w_4(E)=0$ we need to use some dimensional count.
Let $p$ (resp. $q$) denote the ramification point of $D$ contained in $D_1$ (resp. $D_2$). Let $(2p,D_1)$ (resp. $(2q,D_2)$) denote the degree $2$ zero-dimensional subscheme of $D_1$ (resp. $D_2$) with $p$ (resp. $q$) as its reduction. Set $Z:= (2p,D_1)\cup (2q,D_2)$. Since $p$ and $q$ are ramification points of $D$, $Z$ is a degree $4$ subscheme of $D$. Hence there is an exact sequence of sheaves
\begin{equation}\label{eqb1}
0\to \mathcal O_{\mathbb P^1\times \mathbb P^1}\to \mathcal I_Z(n,1)\to \mathcal I_{Z,D}(n,1)\to 0
\end{equation}
 The K\"{u}nneth formula gives $h^1(\mathcal O_{\mathbb P^1\times \mathbb P^1})=0$. Since $\deg(Z)=4$, $D\cong \mathbb P^1$, $n\ge 2$ and $\deg(\mathcal O_D(n,1)) =2n$, we have
$h^1(D,\mathcal I_{Z,D}(n,1)) =0$, i.e. $\dim |\mathcal I_{Z,D}(n,1)| =\dim |\mathcal O_{\mathbb P^1\times \mathbb P^1}(n,1)|-4$. 

Take $p'\in D_1$ and $q'\in D_2$ and set
$Z(p',q'):= (2p',D_1)\cup (2q',D_2)$. By the semicontinuity theorem for cohomology  \cite[Th. III.13.8]{h} we have  $h^1(\mathcal I_{Z(p',q')}(n,1))=0$
for all $(p',q')$ in a Zariski neighborhood of $(p,q)$ in $D_1\times D_2$. Hence there is an irreducible locally closed and codimension $2$ algebraic subset $\Gamma_1$ of $|\mathcal O_{\mathbb P^1\times \mathbb P^1}(n,1)|_{sm}$ such that $D\in \Gamma_1$ and
all $X\in \Gamma_1 $  have an ordinary ramification point contained in $D_1$ and an ordinary ramification point contained in $D_2$. 

Since $D$ has only ordinary ramification points and $p$ and $q$ are its only ramification points contained in $D_1\cup D_2$, there is a Zariski open neighborhood $\Gamma_2$ of $D$ in $\Gamma_1$ such that all $C\in \Gamma_1$ have only ordinary ramification points, $D_1$ contains only one ramification point of $C$ and $D_2$
contains only one ramification point of $C$.

For $j=3,4$ let $\Gamma(j)$ denote the set of all $C\in \Gamma_2$ which are transversal to $D_j$. Since transversality is an open condition for the Zariski topology, $\Gamma(j)$ is Zariski open in $\Gamma_2$. Note that any bundle $E$ with $G(E)\in \Gamma(3)\cap \Gamma(4)$ satisfies $w_1(E)=w_2(E)=1$ and $w_3(E)=w_4(E)=0$. Assume for the moment $\Gamma(3)\ne \emptyset$ and $\Gamma(4)\ne \emptyset$. Since $\Gamma_2$ is irreducible of dimension $\dim |\mathcal O_{\mathbb P^1\times \mathbb P^1}(n,1)|-2$, the same is true for $\Gamma(3)\cap \Gamma(4)$. Hence part (a) holds if $\Gamma(3)\ne \emptyset$ and $\Gamma(4)\ne \emptyset$.
Assume for instance $\Gamma(3)=\emptyset$, i.e. assume that all $C\in \Gamma_2$ are tangent to $D_3$. In particular this holds for $D$. Call $a$ the point of $D_3\cap D$ at which $D_3$ and $D$ are tangent. Let $(2a,D_3)$ denote the degree $2$ zero-dimensional subscheme of $D_3$ with $a$ as its reduction. Set $W:= Z\cup (2a,D_3)$. Since $p$, $q$ and $a$ are ramification points of $D$, $W\subset D$. Hence there is an exact sequence of sheaves
\begin{equation}\label{eqb2}
0\to \mathcal O_{\mathbb P^1\times \mathbb P^1}\to \mathcal I_W(n,1)\to \mathcal I_{W,D}(n,1)\to 0
\end{equation}
 The K\"{u}nneth formula gives $h^1(\mathcal O_{\mathbb P^1\times \mathbb P^1})=0$. Since $\deg(Z)=4$, $D\cong \mathbb P^1$, $n\ge 3$ and $\deg(\mathcal O_D(n,1)) =2n$, we have
$h^1(D,\mathcal I_{W,D}(n,1)) =0$, i.e. $\dim |\mathcal I_{W,D}(n,1)| =\dim |\mathcal O_{\mathbb P^1\times \mathbb P^1}(n,1)|-6$. Since $\dim (D_1\times D_2\times D_3)=3$, varying the points $(p',q',a')
\in D_1\times D_2\times D_3$, we get that $\Gamma_2\setminus \Gamma(3)$ has codimension at least $1$ in $\Gamma_2$.
Hence $\Gamma(3)\ne \emptyset$, contradicting one of our assumptions. \\

 (b) In this step we prove part (b). Let $\gamma_2$ be the only element of  $\mathrm{Aut}(\mathbb P^1)$ such that $\gamma_2(a_1)=\pi_2(D_1)$, $\gamma_2(a_2)= \pi_2(D_2)$ and $\gamma_2(a_3)=\pi_2(D_3)$. Let $\gamma$ be the automorphism of $\mathbb P^1\times \mathbb P^1$ which acts as the identity on the first factor and as $\gamma_2$ on the second factor. Set $X:= \gamma(A)$. 
 
 Let $p$ (resp. $q$, resp $a$) denote the contact locus of $X$ and $D_1$ (resp. $D_2$, resp. $D_3$). Note that $w_1(E)=w_2(E)=w_3(E)=1$ for all bundles $E$ such that $G(E)=X$. Hence to prove part (b) it is sufficient to count the curves $X'$ near $X$ and tangent to $D_j$ for $j=1,2,3$ and prove that the general such curve is not tangent to $D_4$. Let $p$ (resp. $q$, resp. $a$) be the point of contact of $X$ and $D_1$ (resp. $D_2$, resp. $D_3$). As in step (a) let $(2p,D_1)$ (resp. $(2q,D_2)$, resp. $(2a,D_3)$) denote the degree $2$ connected zero-dimensional subscheme of $D_1$ (resp. $D_2$, resp. $D_3$) with $p$ (resp. $q$, resp. $a$) as its reduction. Set 
 $$W:= (2p,D_1)\cup (2q,D_2)\cup(2a,D_3).$$ 
 Since $X$ ramifies at $p$, $q$ and $a$, $W\subset X$. We get an exact sequence similar to \eqref{eqb2} with $X$ instead of $D$.
As in step (a) we get $h^1(\mathcal I_W(n,1)) =0$, i.e. $\dim |\mathcal I_W(n,1)| =\dim |\mathcal O_{\mathbb P^1\times \mathbb P^1}(n,1)|-6$. For any$(p',q',a')\in D_1\times D_2\times D_3$ and set 
$$W(p',q',a'):= (2p',D_1)\cup (2q',D_2)\cup(2a',D_3).$$ By the semicontinuity theorem for cohomology  \cite[Th. III.13.8]{h} we have  $h^1(\mathcal I_{W(p',q',a')}(n,1)) =0$, i.e. $\dim |\mathcal I_{W(p',q',a')}(n,1)| =\dim |\mathcal O_{\mathbb P^1\times \mathbb P^1}(n,1)|-6$ for all $(p',q',a')$ in a non-empty open subset of $(p,q,a)$ in  $D_1\times D_2\times D_3$. Since $\dim (D_1\times D_2\times D_3)=3$, we get
a codimension $3$ analytic subset $\Delta'$ of $\mathcal M_n(0)$ such $w_1(E)=w_2(E)=w_3(E)=1$. 

To prove step (b) we need to prove that, outside a lower dimensional analytic subset of $\Delta'$, the bundles satisfy $w_4(E)=0$, i.e. their graph is transversal to $D_4$. As in step (a) we have a locally closed irreducible subset $\Delta_2$ of $|\mathcal O_{\mathbb P^1\times \mathbb P^1}(n,1)|(0)$ with codimension $3$, containing $X$ and formed by smooth curves with ordinary ramifications, $2n-2$ images of the ramification points. We need to prove that a general $X'\in \Delta_2$ is transversal to $D_4$. Assume that this is not true. In particular $X$ is tangent to $D_4$. Since the ramification points of $X$ have different, there is a unique $b\in D_4\cap X$ at which $X$ and $D_4$ are tangent.
Let $(2b,D_4)$ denote the degree $2$ zero-dimensional subscheme of $D$ with $b$ as its reduction. Set $W':= W\cup (2b,D_4)$.
Note that $\deg(W')=8$ and $W'\subset X$. We get an exact sequence of sheaves similar to \eqref{eqb2} with $X$ instead of $D$ and $W'$ instead of $X$.
Since $n\ge 4$, $X\cong \mathbb P^1$, $\deg(W')=8$ and $\deg(\mathcal O_X(n,1)) =2n$, we have $h^1(X,\mathcal I_{W',X}(n,1)) =0$ and hence
$h^0(\mathcal I_{W'}(n,1)) =h^0(\mathcal I_W(n,1)-2$. Since $\dim D_4 =1$, repeating the last part of step (a) we get a contradiction.
\end{proof}

\begin{thebibliography}{99}
\bibitem[BH]{BH} P. J. Braam, J. Hurtubise, {\it Instantons on Hopf surfaces and monopoles on solid tori}, J. Reine Angew. Math. {\bf 400} (1989) 146-172.
\bibitem[BMo]{BMo} V. Br\^{\i}nz\u{a}nescu, R. Moraru, {\it Stable bundles on non-K\"ahler elliptic surfaces}, Commun. Math. Phys. {\bf 254}  n.3   (2005) 565--580.
\bibitem[H]{h} R. Harshorne, {\it Algebraic Geometry}, Springer, Berlin, 1977.
\end{thebibliography}
\end{document}
